\documentclass[a4paper,12pt,twoside]{article}
\usepackage[left=1in,right=1in,top=1in,bottom=1in]{geometry}

\input{../styles/preamble.tex}
\input{../styles/Chinese_math_style.tex}
\input{../styles/math_style.tex}

\newcommand{\redtext}[1]{{\color{red} #1}}
\newcommand{\bluetext}[1]{{\color{blue} #1}}

\fancypagestyle{firstpage}{
  \fancyhf{}
  \fancyhead[R]{\small《高中數學實驗教材》第 1 冊}
  \fancyhead[L]{\small 第 1 章\quad 引言}
  \fancyfoot[C]{\thepage}
}

\pagestyle{fancy}
\renewcommand{\headrulewidth}{0.5pt}
\renewcommand{\footrulewidth}{0pt}
\fancyhead{}
\fancyhead[LE]{\small 《高中數學實驗教材》第一冊}
\fancyhead[RO]{\small 第一章\quad 引言}
\fancyfoot[C]{\thepage}

\title{\LARGE\textbf{《高中數學實驗教材》第一冊}\\[0.5em] \Large\textbf{第一章\quad 引言}}
\author{\textbf{高中數學實驗教材編輯小組}\quad 編著}
\date{民國六十二年（1973）修訂版}

\begin{document}

\maketitle
\thispagestyle{firstpage}

% ==========================================
% 掃描檔勘誤、文字校訂與疑難處說明（依規則 6 整理於目錄之前）
% ==========================================
\section*{【說明】掃描檔勘誤、疑難處與文義校訂表}
\addcontentsline{toc}{section}{【說明】掃描檔勘誤、疑難處與文義校訂表}
\label{sec:errata_notes}

本章依據 1973 年版《高中數學實驗教材 第一冊》（高一上）第一章「引言」掃描原件轉錄為 \LaTeX{} 格式。為忠實反映教材原始內容並兼顧排版品質與嚴謹性，依作業規範將掃描檔中出現之印刷訛誤、符號缺漏、字模損毀、歷史印記雜訊及文義疑難處詳細考證彙整於下表。讀者可點擊表格中之章節連結直接跳至正文對應位置：

\begin{table}[H]
\centering
\small
\renewcommand{\arraystretch}{1.35}
\begin{tabularx}{\textwidth}{|p{3.2cm}|p{3.6cm}|X|}
\hline
\textbf{章節位置} & \textbf{掃描檔原始現象 / 誤植} & \textbf{考證說明與校訂處理} \\ \hline

\hyperref[sec:1-1]{\S 1-1 頁首} 與 \hyperref[sec:ex1-1]{習題 1-1 頁尾} & 
出現「\texttt{946}」、「\texttt{白單]飯麵}」、「\texttt{货西東}」等破碎字樣 & 
此為原書圖書館藏印、流通戳記、背面透印或早期 OCR 掃描所生之雜訊，正文中已予以剔除，恢復正文字樣。 \\ \hline

\hyperref[subsec:number]{\S 1-1 甲、數} & 
「叫做實系，實數系的性質是數學的重要之基礎之一」 & 
「實系」顯為「實數系」排版之脫字，正文中校補為「實數系」，使數學名詞體系完整。 \\ \hline

\hyperref[fn:sqrt2]{\S 1-1 甲 註釋} & 
「由畢氏定理可知 ZM的長度是AB的2倍...證明/不可能是有理！」 & 
原書根號印損為「2倍」或斜線「/」，且末尾「有理」漏「數」字。依畢氏定理與第二冊代數銜接，校正為「$ZM$ 的長度是 $AB$ 的 $\sqrt{2}$ 倍」及「$\sqrt{2}$ 不可能是有理數」。 \\ \hline

\hyperref[subsec:function]{\S 1-1 乙、函數 (i)-(iv)} & 
「其『體』即隨之確定」、「『體』與『壓力』成反比」、「相互關可以由...」 & 
(i)、(ii) 處將「體積」簡印為「體」；(iv) 處「相互關」脫漏「係」或「連」字（應為「相互關係/相互關連」）。正文均予補全。 \\ \hline

\hyperref[def:function]{函數的定義} 與 \hyperref[subsec:function]{\S 1-1 乙 虎克定律} & 
「且稱y爲應變數\texttt{*}爲自變數」、「重量用:表示」、「從彈簧所填之長算出物重」 & 
「\texttt{*}」及「:」均為變數 $x$ 筆劃印損斷裂；「所填之長」實為「所伸之長」之形近排印訛誤，已予校正。 \\ \hline

\hyperref[subsec:space]{\S 1-1 丙、空間} & 
「弄清楚地到乙地的方向」、「毫無障 硬的平地」 & 
「清楚地」脫漏「甲」字（應為「清楚甲地」）；「障 硬」為鉛字「障礙」破損，正文補正。 \\ \hline

\hyperref[sec:1-2]{\S 1-2 數學方法 引言} & 
「過上面向同學們簡略地介紹了...談一談討這些題材」 & 
開頭「過」字為殘留衍字，「談一談討」為「談一談探討」之排印重複脫漏。 \\ \hline

\hyperref[subsec:logic]{\S 1-2 乙、邏輯} & 
「如此一部有用的機器」、「十之八九你 凑成一具...」 & 
「如此」依前後文意為「好比」之訛；「十之八九你」後脫漏「會」字，正文已潤飾補正。 \\ \hline

\hyperref[subsec:math_language]{\S 1-2 丙、數學語言} & 
「討論抽象數學語言...」、「數學語 雖然...」、「鑽 管」 & 
標點符號脫漏；「數學語」漏「言」字；「鑽 管」為成語「鑽牛角尖」之字模脫落。 \\ \hline

\hyperref[sec:1-3]{\S 1-3} 與 \hyperref[sec:1-4]{\S 1-4} & 
項目編號印為「(C)」；開頭殘印「媽:」；要點 (二)、(三) 漏編號標籤 & 
\S 1-3 第二點校訂為 (二)；清除「媽:」污痕，並重整 \S 1-4 四大要領之編號結構。 \\ \hline

\hyperref[sec:ex1-1]{習題 1-1 第 6 題} & 
「放在盡紙上嗎？」、「把已看成“經過抽象”」 & 
「盡紙」為「畫紙」之形近誤植；「把已」校正為代名詞「把它」。 \\ \hline

\end{tabularx}
\end{table}

\clearpage
\tableofcontents
\newpage

% =========================================================================
% 《高中數學實驗教材》第一冊（高一上） 
% 第一章 引言
% =========================================================================

\section{引言}\label{sec:ch1}

\subsection{數學題材}\label{sec:1-1}

在同學們開始有系統地學習三年高中數學課程之前，開宗明義讓我們先介紹數學的題材，使大家對今後三年數學課程中所要討論的對象——「數」、「函數」、「空間」——有一個概略而整體性的認識。

\subsubsection*{甲、數}\label{subsec:number}

在日常生活中，我們會經常遭遇到種種「量」的問題，譬如你的身高、體重，你家到學校有多遠？家中養的雞今天一共生了幾個蛋？等等。一般說來，常見的「量」可以歸成兩大類：譬如一羣牛，一籃雞蛋，一排樹；它們都具有天然的「個別單位」（即一隻牛，一個蛋、一棵樹）；對於這一類具有天然個別單位的量，我們只需要「數一數」它們有多少個，而用來數（ㄕㄨˇ）個數的工具就是同學們所熟知的「自然數」$1, 2, 3, 4, \dots$。

另一類如「長度」、「重量」、「體積」等，這一類的「量」，不必具有天然而又便於使用的單位，比如說穀子，我們知道它可以個別地以顆粒計算，然而有什麼人這樣計算呢？大概只有我們在寓言裏看到的那段說給國王聽的故事：一隻蝗蟲爬進穀倉啄了一顆穀子……因此，人為的「單位」便應需要而產生了，單位的選取隨地方而不同，比如說到今天在鄉間還是用斗來量穀子，而在城市裏，我們只知道公斤、毫斤、市斤……，不論如何，有了「人為的單位」之後，所要的事是去求所要度量的量和所取定的「單位」間的比值。今以長度為例說明如下：

\begin{figure}[htbp]
\centering
\begin{tikzpicture}[scale=1.1, >=latex, line cap=round, line join=round]
  % (a) 單位直尺 AB
  \draw[thick] (0, 3.4) -- (2, 3.4);
  \draw[thick] (0, 3.25) -- (0, 3.55) node[above] {$A$};
  \draw[thick] (2, 3.25) -- (2, 3.55) node[above] {$B$};
  \node[above] at (1, 3.55) {\small 單位直尺 $AB$};

  % (b) 線段 EF = 3 個單位
  \draw[thick] (0, 2.3) -- (6, 2.3);
  \foreach \x in {0, 2, 4, 6} {
    \draw[thick] (\x, 2.18) -- (\x, 2.42);
  }
  \node[left] at (0, 2.3) {$E$};
  \node[right] at (6, 2.3) {$F$};
  \node[above] at (3, 2.45) {\small 線段 $EF$（$3$ 個單位長）};

  % (c) 線段 GH = 11/6 個單位
  \draw[thick] (0, 1.2) -- (3.667, 1.2);
  \foreach \i in {0, 1, ..., 11} {
    \draw[thin] ({\i*2/6}, 1.12) -- ({\i*2/6}, 1.28);
  }
  \draw[thick] (0, 1.08) -- (0, 1.32) node[left] {$G$};
  \draw[thick] (3.667, 1.08) -- (3.667, 1.32) node[right] {$H$};
  \node[above] at (1.833, 1.32) {\small 線段 $GH$（$\frac{11}{6}$ 個單位長）};

  % (d) 直角三角形 ZNM (圖一)
  \begin{scope}[shift={(7.2, 0.8)}]
    \coordinate (N) at (0,0);
    \coordinate (M) at (2.2,0);
    \coordinate (Z) at (0,2.2);
    \draw[thick] (N) -- (M) node[midway, below] {$NM = AB$} -- (Z) node[midway, above right] {$ZM = \sqrt{2}AB$} -- cycle node[midway, left] {$ZN = AB$};
    % 直角符號
    \draw (0, 0.3) -- (0.3, 0.3) -- (0.3, 0);
    \filldraw (N) circle (1.5pt) node[below left] {$N$};
    \filldraw (M) circle (1.5pt) node[below right] {$M$};
    \filldraw (Z) circle (1.5pt) node[above left] {$Z$};
    \node[below=10pt] at (1.1,0) {\small 直角三角形（不可公度線段）};
  \end{scope}
\end{tikzpicture}
\caption{長度度量與不可公度線段示意圖}\label{fig:measurement}
\end{figure}

設上面直尺 $AB$ 是我們用來度量長度的單位，而下列線段 $EF, GH, LM$ 是幾個我們所要度量的線段。當我們用直尺 $AB$ 去逐段量 $EF$ 時，發現 $EF$ 恰好可以分成三段和 $AB$ 等長的線段連接而成，我們就說線段 $EF$ 的長度是「$3$ 個單位」。當我們再用直尺 $AB$ 去量線段 $GH$ 時，發現 $GH$ 比一段 $AB$ 長，但又比兩段 $AB$ 短。換句話說，若要度量線段 $GH$，我們所取定的「單位」太大了。假如我們將 $AB$ 等分為 $6$ 段，以其中一段 $AB'$ 為新的「直尺」去量 $GH$；則發現 $GH$ 恰好是由 $11$ 段和 $AB'$ 等長的線段連接而成。於是線段 $GH$ 的長度是 $\frac{11}{6}$ 個單位。

也許同學們會問：是不是任何一個給定的線段 $CD$，都是單位長 $AB$ 的幾分之幾？換句話說，是不是任何一個給定線段 $CD$ 的長度都是單位長 $AB$ 的有理數倍？上述問題的答案是否定的。\cref{fig:measurement} 中 $ZN=NM=AB$ 的直角三角形的斜邊 $ZM$ 就是一個否定上述問題的例子\footnote{\label{fn:sqrt2}註：由畢氏定理可知 $ZM$ 的長度是 $AB$ 的 $\sqrt{2}$ 倍，而在第二冊第三章中我們可以證明 $\sqrt{2}$ 不可能是有理數。}。所以由「量」所產生的數系，要比有理數系更廣，叫做\textbf{實數系}。實數系的性質是數學的重要基礎之一，我們在第二冊中將詳加討論，並說明其應用。

\subsubsection*{乙、函數}\label{subsec:function}

當我們進一步去觀察，去研究自然現象時，我們會發現好些「量」與「量」之間，有密切的關連。譬如以 $2$ 克氫氣為例，當我們細心地去做實驗，而且把在各種溫度與壓力下所測得的體積加以記錄比較，就會發現這 $2$ 克氫氣的「溫度」、「壓力」、「體積」這三種量之間，存在著下列有趣的關連：

\begin{enumerate}[(i)]
  \item 當「溫度」、「壓力」給定後，其「體積」即隨之確定。
  \item 當「溫度」固定不變時，「體積」與「壓力」成反比，亦即
  \begin{equation}\label{eq:boyle}
    \text{體積} = \text{常數} \cdot \frac{1}{\text{壓力}} \qquad \text{〔波義耳定律〕}
  \end{equation}
  \item 當「壓力」固定不變時，其「體積」與「溫度」間的關係如下：
  \begin{equation}\label{eq:charles}
    \text{攝氏 } t \text{ 度時的體積} = (\text{攝氏零度時的體積}) \left(1 + \frac{t}{273}\right) \qquad \text{〔查理定律〕}
  \end{equation}
  （溫度增高時，體積也按比例增大。）
  \item 以 $V$ 表示這個氣體在攝氏 $t$ 度、$P$ 個大氣壓下的體積，具體地說，上述三種量之間的相互關連可以由下列式子表出：
  \begin{equation}\label{eq:ideal_gas}
    V \text{ (公升)} = (22.4 \text{ 公升}) \cdot \frac{1}{P} \cdot \left(1 + \frac{t}{273}\right)
  \end{equation}
  （$2$ 克氫氣在攝氏 $0$ 度、$1$ 個“大氣壓”之下的體積是 $22.4$ 公升）
\end{enumerate}

以上只是一個大家所熟知的簡單的例子，這種「量」與「量」之間的相互關連在自然界是非常普遍的。所以在數學中，我們說明「\textbf{函數}」（function）這個基本概念如下：

\begin{definition}[函數的定義]\label{def:function}
設有兩種量 $x, y$，假如當給定的一個 $x$ 時，$y$ 也有一個隨之而確定，我們便說 $y$ 是 $x$ 的函數，平常以符號 
\begin{equation}\label{eq:function_def}
  y = f(x)
\end{equation}
表示之。且稱 $y$ 為\textbf{應變數}，$x$ 為\textbf{自變數}。
\end{definition}

讓我們再舉一個例子來看，大家都知道把物體放在裝有托盤的懸垂彈簧上彈簧會向下伸長，所放物體重量不同，彈簧伸長也不同，而要彈簧伸長某定長度，就要放相當的重物。國中二年級物理科所學的虎克定律就是說明物體的重量與彈簧伸長長度間的關係，把所伸長度用 $y$ 表示，重量用 $x$ 表示，則在彈簧彈性限度內，
\begin{equation}\label{eq:hooke_y}
  y = cx
\end{equation}
$c$ 稱為彈性係數，與所取彈簧材料有關。若把重量看成自變數，長度 $y$ 便是應變數，其間的關係可用符號 $y = f(x) = cx$ 表示，而說 $y$ 是 $x$ 的函數。

\begin{figure}[htbp]
\centering
\begin{minipage}[t]{0.45\textwidth}
\centering
\begin{tikzpicture}[scale=0.9, >=latex]
  % 天花板固定座
  \fill[pattern=north east lines, draw=none] (-1.2, 4.2) rectangle (1.2, 4.45);
  \draw[thick] (-1.2, 4.2) -- (1.2, 4.2);
  
  % 上掛鉤
  \draw[thick] (0, 4.2) -- (0, 3.8);
  
  % 螺旋彈簧
  \draw[thick, domain=0:7.5*pi, samples=200, smooth, variable=\t] 
    plot ({0.3*sin(\t r)}, {3.8 - 0.09*\t});
  
  % 下掛鉤與秤盤
  \coordinate (SpringBottom) at (0, {3.8 - 0.09*7.5*pi});
  \draw[thick] (0, 1.68) -- (0, 1.15);
  \draw[thick] (-0.8, 1.15) -- (0.8, 1.15);
  \draw[thick] (-0.8, 1.15) -- (0, 1.4) -- (0.8, 1.15);
  
  % 砝碼/重物
  \draw[thick, fill=gray!25] (-0.45, 0.45) rectangle (0.45, 1.15);
  \node at (0, 0.8) {\small 重量 $x$};
  
  % 指針
  \draw[->, thick, red] (0, 1.68) -- (1.6, 1.68);
  
  % 刻度尺
  \draw[thick] (1.8, 0.2) -- (1.8, 4.0);
  \foreach \y in {0.5, 1.0, ..., 3.5} {
    \draw (1.7, \y) -- (1.8, \y);
  }
  \foreach \y in {1.0, 2.0, 3.0} {
    \draw (1.6, \y) -- (1.8, \y);
  }
  \node[right] at (1.8, 3.8) {\small 刻度尺};
  \draw[<->, blue, thick] (2.4, 3.8) -- (2.4, 1.68) node[midway, right] {\small 伸長量 $y$};
  \draw[dashed, gray] (0, 3.8) -- (2.5, 3.8) node[right, black] {\footnotesize 原長};
  \draw[dashed, gray] (0, 1.68) -- (2.5, 1.68);
\end{tikzpicture}
\caption{螺旋彈簧的實驗}\label{fig:spring}
\end{minipage}
\hfill
\begin{minipage}[t]{0.5\textwidth}
\centering
\begin{tikzpicture}[scale=0.9, >=latex]
  % 坐標軸
  \draw[->, thick] (-0.3, 0) -- (5.0, 0) node[below] {\small 彈簧增長 $y$ (公分)};
  \draw[->, thick] (0, -0.3) -- (0, 4.5) node[left] {\small 外力 $x$ (公克)};
  \node[below left] at (0,0) {\small $O$};
  
  % 刻度
  \foreach \x in {1, 2, 3, 4} {
    \draw (\x, 0.08) -- (\x, -0.08) node[below] {\footnotesize \x};
  }
  \foreach \y in {1, 2, 3, 4} {
    \draw (0.08, \y) -- (-0.08, \y) node[left] {\footnotesize \y};
  }
  
  % 虎克定律直線
  \draw[thick, blue] (0, 0) -- (4.5, 3.6) node[above right] {\small $x = g(y) = \frac{1}{c}y$};
  
  % 示意點與虛線
  \draw[dashed, gray] (2.0, 0) -- (2.0, 1.6) -- (0, 1.6);
  \filldraw[red] (2.0, 1.6) circle (1.8pt);
  
  \draw[dashed, gray] (3.5, 0) -- (3.5, 2.8) -- (0, 2.8);
  \filldraw[red] (3.5, 2.8) circle (1.8pt);
  
  \node[align=center, fill=blue!5, draw=blue!40, rounded corners=2pt, font=\footnotesize] at (2.4, 3.8) {彈性限度內成正比};
\end{tikzpicture}
\caption{彈性體增長量和外力的關係}\label{fig:elasticity}
\end{minipage}
\end{figure}

從另外一方面說，這樣的函數是怎樣求出來的呢？當然是經由實驗，沒有別的方法。得到以上的關係之後，我們可以利用它來稱物體的重量，就是：從彈簧所伸之長算出物重。這時物重可看成應變數，長度 $y$ 是自變數，可用 $x = g(y)$ 表示（見\cref{fig:elasticity}）。若是同一個彈簧，事實上我們有
\begin{equation}\label{eq:hooke_x}
  x = g(y) = \frac{1}{c} y
\end{equation}
$c$ 就是此彈簧的彈性係數。

更廣泛地，如果一量 $w$ 隨著另外幾個量 $x, y, z$ 的給定而隨之確定，則稱 $w$ 是 $x, y, z$ 的函數，以 
\begin{equation}\label{eq:multivar}
  w = f(x, y, z)
\end{equation}
表示之。譬如上述例子中，體積 $V$ 是溫度 $t$ 與壓力 $P$ 的函數；又當溫度固定時，則體積 $V$ 就簡化成一個 $P$ 的函數；反之，如壓力固定，則體積 $V$ 就簡化成 $t$ 的函數。

在高中數學課程中，我們將討論一些常見而基本的函數，並且利用圖形來描出自變數與應變數的關係，使同學對它們有一個具體的初步認識。至於函數做有系統而深入的討論，將是以後「微積分學」的中心課題。

\subsubsection*{丙、空間}\label{subsec:space}

其次我們要講到「空間」。「空間」就是我們生活的所在，我們活動的範圍。譬如你要從甲地趕到乙地，你得先弄清楚甲地到乙地的方向，然後還要知道甲乙兩地之間的距離。再假定兩地之間是毫無障礙的平地，那麼經驗告訴我們，走直線是最短的。又如當我們要造橋鋪路時，自然得作一番精確的測量與設計。所以不論在生活上、工程上，對於我們所生存的空間是應該仔細了解的。

人類對於空間有系統的研究，在希臘時代就有了輝煌的成就，我們也因而沿用當時的名稱，把研究空間的學問叫做\textbf{幾何學}（geometry）。到了十七世紀初，幾何學的研究方法上，由於法國大數學家笛卡兒（Descartes 1596--1650）的創見而獲得重大的改進，這就是\textbf{解析幾何學}（analytical geometry）。解析幾何學的基本原理在於充分利用代數學的方法與精神去研究幾何學。在本書的第三、第四冊中，我們將詳細討論解析幾何學基本原理的精要所在，務求同學對這一數學史上的偉大創見，能有一個深切的認識與領會。

\subsection{數學方法}\label{sec:1-2}

上面向同學們簡略地介紹了數學的基本題材——數、函數與空間，其次我們要和大家談一談探討這些題材的方法，使同學們先有一個扼要的認識；然後在今後討論、學習的過程中，逐步去體會，去學習這些數學方法。

\subsubsection*{甲、抽象}\label{subsec:abstraction}

首先讓我們來談一談什麼叫做「抽象」以及「抽象」在數學中所扮演的角色。當我們討論某一給定的問題時，它往往只跟「該問題的對象」的某幾個「性質」有關，而和該對象同時所具備的其他性質全然無關。那麼，在我們研討這一給定的問題時，把我們的注意力集中在那些有關的幾個性質，而暫時不去理會那些無關的性質是很自然的。把某些事物的重要性質抽提出來，加以深入的探討，從而乾淨俐落地解決有關的許多問題，這就是數學中所常用的\textbf{適度抽象化}的方法。

所以「抽象化」是數學探討中的一種手段。數學的中心課題依然是去瞭解如數、函數、空間等等實在而有用的對象！對於不同的問題的探討，我們作不同程度的（亦即適度的）抽象，那只是為了使我們的討論較為便捷，為了使我們研討時不致被無關的資料迷亂了我們的視界。

今以地圖為比喻，簡易說明「適度抽象化」的要點如下：假如你是初到臺北市的一個陌生人，你希望從臺北車站出發，步行到幾個地方。那麼最簡單的方法是買一張「臺北市街道圖」，因為它就是製圖者把你所要解決的問題的有關資料——通路和幾處要去的地方間的相對位置——抽提出來，象徵性地記述在一張紙上。所以它使你一目瞭然，從而有效地計劃一個最簡捷的走法。但是假如你要去的地方較遠，你希望乘公共汽車去，那麼在你買臺北市地圖時就得注意它是否在圖上同時記載著公共汽車的路線。最後，假如你是位越洋的飛行員，你所需要的是一幅可資導航用的世界地圖；一張大街小巷全畫得清清楚楚的世界地圖不但大得無法攜帶，而你假如真的帶了這麼一張極笨重的世界地圖的話，你的飛行速度也將使你閱圖時疲於奔命呢！

從上述的比喻，同學們應該大略地體會到「抽象化」的「適度性」是十分重要的。至於在數學討論上的例子，留待適當時機時，我們再向同學們介紹。

\subsubsection*{乙、邏輯}\label{subsec:logic}

其次讓我們來談談「邏輯」在數學中所扮演的角色。「邏輯」兩字，是 logic 的譯音，而邏輯學乃是討論思考法則的學問。在一般數學或其他學科中，簡易的邏輯（請參看 \hyperref[sec:2-3]{\S 2-3}）是一個基本的工具。同學並不需要先熟讀「邏輯學」再來學習數學或其他學科如物理、化學、經濟學等等。但是對於邏輯的最簡易基本部分卻一定要弄得清清楚楚，並逐漸熟習其用法。

「簡易邏輯」在數學中所佔的角色大致說來是這樣的：譬如說在幾何學中我們研究空間的性質，但是空間的種種性質何止千萬？而且不斷地研究還會不斷地發現新的性質。假如我們不設法加以整理，只是把空間的無數性質拉拉雜雜地收集著，最後勢必成了猴子撿果子，隨撿隨丟。所以我們必須要把空間的種種性質加以整理，而\textbf{邏輯推理}（deduction）就是我們的工具，我們的不二法門。這種整理法在古希臘時代，即已被廣泛應用，而歐幾里得（Euclid）利用邏輯推理對於當時的幾何知識所作的整理，就成為數學史上一部曠古巨著——《幾何原理》（\textit{Elements}）。數千年來，《幾何原理》中所用的方法一直是數學中治學的經典範本，簡略解釋如下：

假如性質 $M$ 是另外幾個性質 $A, B, C$ 的必然結論；換句話說，當任一事物同時具有性質 $A, B, C$ 時，該事物就必然也具有性質 $M$，則我們說 $M$ 是 $A, B, C$ 的一個推論（或者說 $M$ 可以從 $A, B, C$ 由邏輯推演而得）。理論上，假如我們充分瞭解性質 $A, B, C$，那麼我們也因而知道了基於性質 $A, B, C$ 所能推演得的種種性質。所以，很自然的，我們會問：「是不是可以整理出一組空間的性質 $A, B, C, \dots$ 等，使得其他空間的性質都是上述這組性質的邏輯推論？」這也就是歐氏很成功地在《幾何原理》一書中所完成的工作。（他當時把所整理出的一組基本性質叫做幾何公理，其實還不如直截了當叫做「基本性質」來得妥當些。）

現在且用下述比喻，粗淺地說明「邏輯推理」在數學中的角色：數學的目的是要在瞭解前述題材如數、函數、空間等等的有用性質。一個有用的性質（或定理）好比一部有用的機器，而那些整理出來的基本性質就好比構成機器所需的各種零件，邏輯推理便好比適當地連接那些零件的螺絲釘。一個造機器的人一定得會安裝螺絲釘。數學中有效而廣泛地運用簡易邏輯，但是數學絕非邏輯遊戲，就好像你假如盲目地把一堆零件用螺絲釘去瞎拼，十之八九你會湊成一具無用而笨重的怪物！

\subsubsection*{丙、數學語言}\label{subsec:math_language}

語言文字是溝通思想的橋樑與表達意思的工具，有語言才能講述邏輯，才能討論抽象。數學語言整個看來具有日常語言的形式與本質，但局部地，卻常有它特別的用法，就像法律條文是在闡述法律，醫學名詞是用來講析病理，樂譜是用來記述歌調樂章一樣。數學語言雖然已獨自發展成一門學問，但是我們不需要也不應該在這方面鑽牛角尖。我們固然要熟識所需要的符號，但千萬不要因未習慣這些符號而發生困難時就以為「侯門深似海」，認為自己永無登堂入室的份了。

事實上，法律條文並不是法律，不懂得法律條文的人也有充分權利受到法律的保障，但我們還是應該儘量充實個人的法律常識；醫學名詞固然是醫生的專長，但絕不應該是醫生用來聚斂財產的手段，以及用來炫耀其學識的工具，正相反，它可以用來向大眾闡明健康和衛生的知識的，現代人的醫學常識不是比一百年前的人多得多嗎？同樣的，我們看待數學語言要注重的是它所敘述的內容與推理的根據，亦即它所根源的自然與人文的現象與所帶引出來的問題。

「數學方法」並非僅限於適度抽象化與邏輯推理，但是這兩者是其中最基本、最簡單也是最有用的。所以同學們務必對上述二者隨時注意，多下功夫學習，去體會如何活用它們。（其他的數學方法，留待真正用到時，再就近取材為同學們詳加說明。）求學一如行遠，治學的方法一如步法，只要切實學會穩健的步法，有毅力地堅定走下去，看準方向，一定會達到目的地的。

\subsection{中學數學教育的目的}\label{sec:1-3}

現在讓我們談一談為什麼高中三年，每學期都安排著很重的數學課程，其目的何在？

\begin{enumerate}[(一)]
  \item 數學中的基本題材如數、函數、空間是極基本，極有用的。不論日常的問題或者是工程上、自然科學上、社會科學上的問題，它們都是不可或缺的有效工具，高中數學教育的一個明顯的目的是確保同學們對上述基本題材的一個有系統的明確瞭解，並從而熟悉其應用。
  \item 另外一個常常被同學們所忽略的重要目的，便是思考能力的訓練，因為數學中的題材簡單明確，而其中涉及「推理」、「分析」、「綜合」……等思考方法時，每每能清晰地顯示其用法與精微之處。故數學一直是中學課程中闡述思考方法，訓練同學們獨立思考能力的主要科目。
\end{enumerate}

\subsection{學習數學應有的態度及應把握的要點}\label{sec:1-4}

在六年的中學教育中，每年都有繁重的數學課程，而對許多同學來說，數學是叫人頭痛的科目。加上它偏偏又是一門主科，於是好些同學對數學又厭又懼，卻又不得不天天下功夫、傷腦筋。但是往往由於不得其法，愈學愈糊塗，臨時只好去死記硬背些題目，希望僥倖圖個及格，這樣六年下來，中學教育豈非一場噩夢！

同學們！讓我們從今天開始，對學習數學抱定一個健全、正確的態度：\textbf{數學的題材是實在的，有用的；它們的基本性質簡單，完美；數學的理論是逐步漸進的、條理分明的}。所以只要同學們肯下功夫去練習，數學是一門有趣、有用而且易學的科目！當然下功夫還是要得法，要把握下列幾點要領：

\begin{enumerate}[(一)]
  \item 科學的基本精神在於以「簡」御「繁」，以「已知」去瞭解「未知」。所以在學習時，基本而簡單的部分是最重要的，務求徹底弄懂，唯有把「簡單的」弄得一清二楚，才能用來解決「繁難的」。許多同學對數學感到困難的一個原因就是他們忽略了基本定義與基本性質，未能好好下功夫弄清楚。因此到較深的部分就成了一團糟的局面。所以當同學在學習某一段感到困難時，首先得把其中用到的概念與事實再弄弄清楚，然後再下功夫反覆想想，也就不難弄明白了。
  \item 由於數學的定理是逐步漸進的，大多數使同學發生問題的緣故，都不是現在所講的東西太難、太深，而是過去所學的東西，沒有徹底明瞭，解決這項困難的唯一辦法，是在不清楚時，隨時複習過去的資料。
  \item 前面已經講過，數學對於問題的處理，要適度地加以抽象化。但是數學所研討的問題，大多數有他們的來源與直觀涵義，同學們對此千萬不可忽略。
  \item 做習題是幫助了解、訓練技巧所不可缺少的工作，做習題時，必須注意解題的思路與一般法則，才能收到舉一反三的效果。習題一定要多做，多練習，才能抓到要點，而不是死背習題與例題，那只說明他不能理解解題的精要所在，而且，當情況改變時，他便不一定能解題了。
\end{enumerate}

除上述幾點之外，數學跟任何其他學問一樣，學習時應該多問多想，與良師益友相切磋，更能收事半功倍之效。

\subsection*{習題 1-1}\label{sec:ex1-1}
\addcontentsline{toc}{subsection}{習題 1-1}

\begin{enumerate}[1., leftmargin=2em]
  \item 想一想，學習的興趣是否與考試的分數有關？\label{prob:1}
  \item 想一想，考試的分數是否與準備的充分有關？\label{prob:2}
  \item 當你唸到一段感到困難，只知道不懂，可是不知道哪一點不懂時，就是不知道問題在哪裏，那時怎麼辦？\label{prob:3}
  \item 如果你連著考試不及格，你便要認定自己笨、智商低，因此灰心不想再唸下去，是必然的嗎？有沒有別的出路？\label{prob:4}
  \item 如果所講的題材在國中時也讀過，應該怎麼對待、處理？\label{prob:5}
  \item 以你寫生的經驗來說，你能如照相一樣把所見的全都放在畫紙上嗎？你畫的不是只有幾個線段幾抹形象？這些不論能不能構成一幅好的畫面，是不是都可以把它看成“經過抽象”而得的？\label{prob:6}
  \item 想一想，數學經過「抽象」後，是否可以說明實際情形？（注意：如\cref{prob:6}，這裡說明抽象不是「美化」的工作，而且是抓重點的工作。）\label{prob:7}
\end{enumerate}


\end{document}
